%%%%%%%%%%  scan idx 297  |  folio 284  %%%%%%%%%%
\SGAPrintPage{284}%
\((1=\cl(\underline{O}_X))\), \uline{and the canonical homomorphism} (2.6.3.2)
\[
 K^{\bullet}(X)_{\naif}\longrightarrow K^{\bullet}(X)
\]
\uline{is a homomorphism of unital rings}.

Let \(f:(X,\underline{O}_X)\to(Y,\underline{O}_Y)\) be a morphism of topoi ringed by commutative rings. Then compatibility of the tensor product with inverse images shows that \uline{the homomorphism}
\[
 f^*:K^{\bullet}(Y)\longrightarrow K^{\bullet}(X)
\]
(resp.
\[
 f^*:K^{\bullet}(Y)_{\naif}\longrightarrow K^{\bullet}(X)_{\naif}
\]
) \uline{is a homomorphism of unital rings}.
\end{enumerate}

\sgaStableTarget{sga6.e04.para.2-8}2.8. \uline{Augmentation of} \(K^{\bullet}\) \uline{to} \(H^0(\ ,\mathbb Z)\). Let \((X,\underline{O}_X)\) be a locally ringed topos (I 2.15.1). Recall that we have defined (I 6.11 to 6.14) a homomorphism of (unital) rings
\[
 i:H^0(X,\mathbb Z)\longrightarrow K^{\bullet}(X)_{\naif},
\]
thus giving \(K^{\bullet}(X)_{\naif}\) (and consequently also \(K^{\bullet}(X)\) via the canonical homomorphism \(K^{\bullet}(X)_{\naif}\to K^{\bullet}(X)\)) the structure of an algebra (with unit) over \(H^0(X,\mathbb Z)\). Moreover (loc.cit.), the algebra \(K^{\bullet}(X)_{\naif}\) (resp. \(K^{\bullet}(X)\)) is augmented to \(H^0(X,\mathbb Z)\) by the \uline{rank}, and the canonical homomorphism \(K^{\bullet}(X)_{\naif}\to K^{\bullet}(X)\) is a homomorphism of augmented algebras.

Let \(f:(X,\underline{O}_X)\to(Y,\underline{O}_Y)\) be a morphism of locally ringed topoi. The homomorphisms \(f^*:K^{\bullet}(Y)\longrightarrow K^{\bullet}(X)\) and \(f^*:K^{\bullet}(Y)_{\naif}\longrightarrow K^{\bullet}(X)_{\naif}\) are compatible (I 6.12, 6.13) with the preceding augmented algebra structures; this is expressed by the commutativity of the diagram (where \(\Rg\) denotes rank):

%%%%%%%%%%  scan idx 298  |  folio 285  %%%%%%%%%%
\SGAPrintPage{285}%
\[
\begin{tikzcd}[column sep=tiny]
H^0(Y,\mathbb Z)\arrow[r,"i"]\arrow[d]
&K^{\bullet}(Y)_{\naif}\arrow[r]\arrow[d]
&K^{\bullet}(Y)\arrow[r,"\Rg"]\arrow[d]
&H^0(Y,\mathbb Z)\arrow[d]\\
H^0(X,\mathbb Z)\arrow[r,"i"]
&K^{\bullet}(X)_{\naif}\arrow[r]
&K^{\bullet}(X)\arrow[r,"\Rg"]
&H^0(X,\mathbb Z).
\end{tikzcd}
\]

\textsuperscript{*}We shall see later in the Seminar (V 3.9.1 and VI) that \(K^{\bullet}(X)_{\naif}\) can be given a richer structure than that of an algebra augmented to \(H^0(X,\mathbb Z)\), namely that of a \(\lambda\)-algebra augmented to \(H^0(X,\mathbb Z)\) regarded as a binomial ring (V 2.8), this structure being compatible with inverse-image homomorphisms. Deligne [1] recently proved that \(K^{\bullet}(X)\) can be given an analogous structure (at least when \(X\) has characteristic \(p\geq0\)) and that the canonical homomorphism \(K^{\bullet}(X)_{\naif}\longrightarrow K^{\bullet}(X)\) is a homomorphism of augmented \(\lambda\)-algebras, thus answering a conjecture of Grothendieck (XIV 1) in the affirmative.\textsubscript{*}

\sgaStableTarget{sga6.e04.para.2-9}2.9. \uline{Comparison of} \(K^{\bullet}\) \uline{and} \(K^{\bullet}_{\naif}\). Let \((X,A)\) be a ringed topos. In general there is no reason for the canonical homomorphism
\[
 K^{\bullet}(X)_{\naif}\longrightarrow K^{\bullet}(X)
\]
to be injective or surjective. However, the ``criteria for the existence of global resolutions'' of Exposé II show that it is an isomorphism in the following cases:
\begin{enumerate}[label=(\roman*)]
\item \(X\) is the Zariski topos of a divisorial scheme, separated or noetherian (and \(A=\underline{O}_X\)) (II 2.2.9);
\item \(X\) is the topos defined by a compact space and \(A\) is a soft sheaf (II 2.3.2);
\item \(X\) is a Stein topos whose final object is quasi-compact and \(A=\underline{O}_X\) (II 2.4.5).
\end{enumerate}

%%%%%%%%%%  scan idx 299  |  folio 286  %%%%%%%%%%
\SGAPrintPage{286}%
\sgaStableTarget{sga6.e04.para.2-10}2.10. \uline{Pairing between} \(K^{\bullet}\) \uline{and} \(K_{\bullet}\). Let \((X,\underline{O}_X)\) be a commutatively ringed topos. We know (I 5.8.4 b)) that there is a pairing
\[
 \otimes_A^{\mathrm L}:D(X)_{\uparf}\times D^b(X)_{\coh}
 \longrightarrow D^b(X)_{\coh}.
\]
This defines a homomorphism
\sgaStableTarget{sga6.e04.eq.2-10-1}\begin{equation}\tag{2.10.1}
 K^{\bullet}(X)\otimes K_{\bullet}(X)\longrightarrow K_{\bullet}(X).
\end{equation}
Thanks to the associativity of the tensor product, we see that (2.10.1) \uline{endows} \(K_{\bullet}(X)\) \uline{with the structure of a} \(K^{\bullet}(X)\)-\uline{module}. More generally, if \(f:(X,\underline{O}_X)\longrightarrow(Y,\underline{O}_Y)\) is a morphism of commutatively ringed topoi, we can regard \(K_{\bullet}(X)\) as a \(K^{\bullet}(Y)\)-module via the homomorphism \(f^*:K^{\bullet}(Y)\longrightarrow K^{\bullet}(X)\).

\sgaStableTarget{sga6.e04.para.2-11}2.11. \uline{Covariant law on} \(K_{\bullet}\) \uline{and projection formula}. Let \((X,A)\) be a ringed topos. Unlike \(K^{\bullet}({}_A X)\), the group \(K_\bullet({}_A X)\) has no natural variance with respect to \((X,A)\). However, in Algebraic Geometry and Analytic Geometry, \(K_{\bullet}(X)\) tends to be \uline{covariant} in \(X\) for \uline{proper} morphisms. Let us make this precise:

\sgaStableTarget{sga6.e04.para.2-11-1}2.11.1. Let \(f:X\to Y\) be a proper and pseudo-coherent morphism (III 1.2) of quasi-compact schemes. Then, subject to the validity of the finiteness conjecture for \(f\) (III 2.1), hence in particular (III 2.2) when \(Y\) is noetherian (in which case the pseudo-coherence hypothesis on \(f\) is automatically satisfied (III 1.3)) or when \(f\) is projective, the functor \(Rf_*\) induces, by (III 2.6), a functor
\[
 Rf_*:D^b(X)_{\coh}\longrightarrow D^b(Y)_{\coh}.
\]
This defines a homomorphism
\sgaStableTarget{sga6.e04.eq.2-11-1-1}\begin{equation}\tag{2.11.1.1}
 f_*:K_{\bullet}(X)\longrightarrow K_{\bullet}(Y)
\end{equation}

%%%%%%%%%%  scan idx 300  |  folio 287  %%%%%%%%%%
\SGAPrintPage{287}%
This homomorphism is \(K^{\bullet}(Y)\)-linear (cf. (2.10)); in other words, for \(y\in K^{\bullet}(Y)\) and \(x\in K_{\bullet}(X)\), we have the so-called \uline{projection formula}
\sgaStableTarget{sga6.e04.eq.2-11-1-2}\begin{equation}\tag{2.11.1.2}
 f_*(yx)=y f_*(x).
\end{equation}
Indeed, if \(x=\cl(E)\) for \(E\in\Ob D^b(X)_{\coh}\) and \(y=\cl(M)\) for \(M\in\Ob D(Y)_{\uparf}\), we have (III 3.7)
\[
 Rf_*(E)\otimes^{\mathrm L}M\simeq
 Rf_*(E\otimes_Y^{\mathrm L}M).
\]
Finally, if \(g:Y\to Z\) is a morphism satisfying the same hypotheses as \(f\) and for which the finiteness conjecture (III 2.1) holds, the transitivity of direct-image functors gives, for the homomorphisms (2.11.1.1), the formula
\sgaStableTarget{sga6.e04.eq.2-11-1-3}\begin{equation}\tag{2.11.1.3}
 g_*f_*=(gf)_*.
\end{equation}

\sgaStableTarget{sga6.e04.para.2-11-2}2.11.2. Let \(f:X\to Y\) be a proper morphism of finite-dimensional analytic spaces. Then, by Grauert's finiteness theorem [2], \(Rf_*\) sends \(D^b(X)_{\coh}\) to \(D^b(Y)_{\coh}\), and therefore defines a homomorphism
\sgaStableTarget{sga6.e04.eq.2-11-2-1}\begin{equation}\tag{2.11.2.1}
 f_*:K_{\bullet}(X)\longrightarrow K_{\bullet}(Y).
\end{equation}
Like its algebraic analogue (2.11.1.1), this is \(K^{\bullet}(Y)\)-linear, i.e. formula (2.11.1.2) is valid, as follows from the projection formula at the level of derived categories (III 5.8 b)). We likewise have the transitivity formula (2.11.1.3).

\sgaStableTarget{sga6.e04.para.2-11-3}2.11.3. Once we leave the algebraic or analytic setting, this covariance of \(K_{\bullet}\) for proper morphisms disappears. Thus, let \(f:X\to Y\) be a proper morphism of topological spaces ringed by the sheaf of complex continuous functions. Then, for \(E\in\Ob D^b(X)_{\coh}\), in general we do not have \(Rf_*(E)\in\Ob D^b(Y)_{\coh}\), as is already evident for \(E=\underline{O}_X\) and \(Y\) reduced to a point. In fact, definition (2.2) of \(K_{\bullet}\)

%%%%%%%%%%  scan idx 301  |  folio 288  %%%%%%%%%%
\SGAPrintPage{288}%
appears especially suited to the needs of Algebraic Geometry and Analytic Geometry, and a good definition of the \(K_{\bullet}\) of a topological space remains to be found. It seems moreover that we shall have to leave the setting of \(\underline{O}_X\)-linear complexes and define, for a morphism \(f:X\to Y\) (of topological spaces, resp. of mixed varieties (II App. II), etc.), a suitable class of complexes of topological vector sheaves on \(X\), which should include, for example, when \(f\) is the structural morphism of a mixed variety, the elliptic complexes on \(X\) relative to \(Y\). This is suggested by the finiteness theorem (II App. II), as well as by a similar theorem (where \(Y\) is supposed reduced to a point) proved by Grothendieck in [3].

\sgaStableTarget{sga6.e04.para.2-12}2.12. \uline{Unusual variances of} \(K^{\bullet}\) \uline{and} \(K_{\bullet}\). Let \(f:(X,A)\longrightarrow(Y,B)\) be a morphism of finite tor-dimension of ringed topoi (III 3.1). The functor \(Lf^*:D^-({}_B Y)\longrightarrow D^-({}_A X)\) sends \(D^b({}_B Y)\) to \(D^b({}_A X)\), hence (I 2.16.1) \(D^b({}_B Y)_{\coh}\) to \(D^b({}_A X)_{\coh}\), and therefore defines a homomorphism
\sgaStableTarget{sga6.e04.eq.2-12-1}\begin{equation}\tag{2.12.1}
 f^*:K_\bullet({}_B Y)\longrightarrow K_\bullet({}_A X).
\end{equation}
If \(A\) and \(B\) are commutative, commutation of inverse image with the tensor product shows that the homomorphism (2.12.1) is compatible with the module structures of \(K_{\bullet}(Y)\) and \(K_{\bullet}(X)\) over \(K^{\bullet}(Y)\) and \(K^{\bullet}(X)\) respectively (2.10); in other words, for \(a\in K^{\bullet}(Y)\) and \(y\in K_{\bullet}(Y)\), we have the formula
\sgaStableTarget{sga6.e04.eq.2-12-2}\begin{equation}\tag{2.12.2}
 f^*(ay)=a f^*(y).
\end{equation}
If \(g:(Y,B)\to(Z,C)\) is a morphism of finite tor-dimension of ringed topoi, transitivity of inverse-image functors implies that the homomorphisms (2.12.1) satisfy the formula:
\sgaStableTarget{sga6.e04.eq.2-12-1-1}\begin{equation}\tag{2.12.1.1}
 f^*g^*=(gf)^*.
\end{equation}

